\section{Project Synthesis: The Judgment of This Study}

\subsection{The judgment}

\begin{judgmentbox}{Project synthesis --- the governing judgment of
this study, with its reasons}
\textbf{This study judges that at Champion the ecclesial judgment and
the documentary record are two different objects, and that the honest
description of the case is: a valid, complete, and unusually candid
diocesan approval, resting deliberately on coherence and fruits rather
than on documents, of an event stratum for which no contemporary
documentation exists and for which the earliest published narratives,
twelve and fifteen years later, disagree with each other and contain
neither the self-identification nor the warning that now stand at the
centre of the received message.} Both halves of that sentence are
load-bearing. Neither half is a concession to the other.\jsep
\textbf{Reason one: the decree says so itself.} It records that ``the
documents from the early history of the Shrine are not abundant''
and answers the objection by appeal to coherence between event,
consequences, character, and mission --- the 1978 Norms' own method of
``cumulative'' and ``convergent'' application of indicative criteria.
The bishop did not claim a documentary case and this study does not
have to impute one to him.\jsep
\textbf{Reason two: the two nineteenth-century witnesses cannot be
harmonized.} Starr has a summer day, eleven miles, an encounter on the
return, a second going and a third returning some weeks later, and a
veiled figure. Pernin has three successive Sundays all on the return,
seven and a half miles, and a date ten or twelve years before 1874.
Neither matches the other; neither matches the received sequence; and
Pernin's date does not even match the received year. A tradition with a
firm early record does not do this.\jsep
\textbf{Reason three: the received message is a composite whose core is
printed English.} The catechetical commission reproduces Starr's 1871
wording almost verbatim; the closing clause reproduces Pernin's 1874
wording; and the self-identification as Queen of Heaven, the general
confession, the punishment clause, the beatitude, and the vineyard
rebuke were not located in any witness this study reached before the
current custodial narrative. Whatever the transmission history, the
words now recited are continuous with print, not independent of it.\jsep
\textbf{Reason four: the fire tradition has a hard core and a soft
shell, and they are separable.} The preservation of the enclosure is
early, circumstantial, and attested by a priest who could be
contradicted by living participants. The rain, the crowd, the
livestock, and the well are not in that witness at all, and a paragraph
supplying the rain circulates under his name that is not his
text.\jsep
\textbf{Reason five: none of this touches the decree's object.} The
decree judges ``the substance'' of the October 1859 events. It adopts
no day, no count, no sequence, no description, and no wording; it
declares nothing a miracle; and it obliges no one. Every finding above
is compatible with it, because the decree was drafted --- carefully, as
its qualifiers show --- so that findings of this kind would be
compatible with it.
\end{judgmentbox}

\subsection{The counterargument at its strongest}

The best case against this study's judgment is not that the collation
is wrong. It is that the collation is beside the point, and it runs
like this.

Adele Brise lived until 1896. She told this story for thirty-seven
years after Starr printed four pages of it in a book for girls, in a
settlement where the chapel, the school, the community, and the annual
processions were all built around the telling. The people who
eventually wrote the fuller version down --- the sisters at the chapel,
the Franciscans who took over in 1902, the historian of 1955 --- stood
in an unbroken chain of custody from the seer herself, in a community
that had heard her tell it hundreds of times. To weigh a Chicago
lecturer's compression of a single interpreted conversation, or a
Montreal fundraiser's self-declared abridgement of ``a long
conversation,'' \emph{against} that chain is to mistake the accident of
publication for the substance of transmission. On this account the
fuller message is not a later accretion at all: it is what Adele
always said, minus what two outside authors, for their own reasons,
left out. The verbal identity with Starr proves only that somebody
later reached for the printed English when writing the story down ---
which is what anyone would do --- not that the content is late. And
Pernin's ``ten or twelve years'' is exactly the kind of error a man
makes about somebody else's biography.

That argument is strong. It cannot be refuted from the sources this
study reached, and it is why every claim in \S5 is phrased as a claim
about the reach of searches rather than about the history of the
tradition. It has one weakness, which this study thinks decisive but
not conclusive: it explains why the fuller material might be old, and
it does not produce a single dated witness that it is. Chains of
custody are arguments about plausibility; they are not evidence about
dates. Until a dated witness appears, the honest description of the
self-identification, the general confession, and the punishment clause
is ``first located in the modern custodial narrative,'' not ``handed
down from 1859.''

\subsection{What would change the judgment}

Four specific findings would move this study, and it is worth naming
them precisely enough that anyone with access to the archives could
check.

\begin{enumerate}
\item \textbf{Any dated document before 1955 containing the
self-identification, the general confession, or the punishment
clause.} A letter, a parish or diocesan record, a devotional pamphlet,
a newspaper notice, a chapel inscription, a pilgrim's account. One such
witness would move those elements out of the ``not located'' row of the
stratigraphy in \S5 and require this study to withdraw its dating.
\item \textbf{Any dated document before 1955 reporting the refugee
crowd, the livestock, the well, or the rain on the night of 8 October
1871.} The same consequence for \S6.
\item \textbf{The source of the paragraph attributed to Pernin.} If it
turns out to be a genuine Pernin text --- a later edition, a periodical
article, a manuscript --- the misattribution finding in \S6 collapses
and this study was wrong. If it turns out to be a twentieth-century
retelling, the finding is confirmed and the attribution should be
corrected wherever it appears.
\item \textbf{Any act of the Holy See concerning Champion.} A record of
Roman authorization or review of the 2010 decree, or any later act
naming Champion, would resolve the question left open in \S7 and would
bear on the standing discussed in \S8.
\end{enumerate}

\subsection{Reading the message within the faith}

The decree found the accounts ``free from doctrinal error and
consistent with the Catholic faith,'' and this study has no basis to
differ. But a message can be free from error and still be misread, and
two of its clauses are misread often enough to deserve a paragraph
each.

\textbf{``I am the Queen of Heaven who prays for the conversion of
sinners.''} Whatever its date, this clause is orthodox exactly as it
stands, and it is orthodox because of what it says Mary does: she
prays. The Council states the boundary: ``There is but one Mediator
\dots{} The maternal duty of Mary toward men in no wise obscures or
diminishes this unique mediation of Christ, but rather shows His
power,'' for her salvific influence ``flows forth from the
superabundance of the merits of Christ, rests on His mediation, depends
entirely on it and draws all its power from it.''\endnote{Second
Vatican Council, \work{Lumen gentium} (21 November 1964), n.~60,
English text read 2026-07-25 at
\url{https://www.vatican.va/archive/hist_councils/ii_vatican_council/documents/vat-ii_const_19641121_lumen-gentium_en.html}.}
A queen who prays is a creature interceding, not a power negotiating.
The message's own logic makes the point: she asks Adele to do what she
does --- ``I wish you to do the same'' --- which is intelligible only
if what she does is something a Christian can imitate.

\textbf{``If they do not convert and do penance, my Son will be obliged
to punish them.''} This is the clause that, read carelessly, produces
the picture the Church has always refused: a merciful mother
restraining an angry Son. Three things block that reading. The clause
is conditional, and conditional warnings of this shape are the ordinary
grammar of scriptural preaching --- the summons to repentance is not a
schedule of disasters. The Son is described as \emph{obliged}, which
locates the necessity in justice and in the sinner's own refusal, not
in a divine temper that intercession must cool. And the mediation is
the wrong way round for the fearful reading: it is the Son's merits
that give the mother's prayer whatever power it has, so that her
intercession cannot be an alternative to him. Where the Council speaks
of Mary as Advocate and Auxiliatrix it immediately adds that this ``is
to be so understood that it neither takes away from nor adds anything
to the dignity and efficaciousness of Christ the one
Mediator.''\endnote{\work{Lumen gentium}, n.~62. The Council adds that
``no creature could ever be counted as equal with the Incarnate Word
and Redeemer,'' and that the Church professes ``this subordinate role
of Mary'' so that the faithful ``may the more intimately adhere to the
Mediator and Redeemer.''} A Champion devotion that leaves a pilgrim
more afraid of Christ than before has inverted its own message.

There is also a reading of the fire that the sources forbid and that
this study will not supply: the fires of 1871 were not a punishment
foretold in 1859. No witness this study reached makes that connection.
Pernin, who did read the catastrophe as a summons to repentance, never
connects it to the apparition, and he had every opportunity: the two
accounts stand within five pages of each other in his own book, and he
draws no line between them.

\subsection{What is left, and it is not small}

Strip Champion of everything this study has been unable to date and
what remains is this. In or about October 1859, in a Walloon
settlement with no priest, an illiterate young Belgian immigrant woman
became convinced that she had been told to gather the children
of the woods and teach them their catechism, the sign of the cross, and
how to approach the sacraments. She did it. She walked between
settlements through rain and snow, was laughed at, taught for free,
begged in Chicago in a black cape-bonnet for a schoolhouse, gathered
five or six women who stayed, built a chapel on the spot, gave the land
to her bishop, and kept it up until she died in 1896. Two independent
outsiders --- one writing for schoolgirls, one for fire relief ---
found the work impressive enough to publish, and both reported the same
commission in substantially the same terms.

That is the part of the record that is not in doubt, and it is exactly
the part the 2010 decree leans on when it speaks of ``coherence between
event and consequences, personality of the seer and commitment to the
mission received.'' It is also, in the end, the shape the message
itself has: not a secret, not a chronology, not a prophecy, but a
commission to teach children what they should know for salvation. A
private revelation that adds nothing to Revelation and sends the
hearer back to the catechism, the sign of the cross, and the sacraments
is doing precisely what the Church says private revelation is for ---
``a help which is proffered,'' whose whole business is ``nourishing
faith, hope and love, which are for everyone the permanent path of
salvation.''\endnote{Benedict~XVI, \work{Verbum Domini} (30 September
2010), n.~14, as quoted at note~24 of the DDF \work{Norms} of 17 May
2024.}
